<<<<<<< HEAD
<<<<<<< HEAD
\documentclass[a4paper, 12pt]{report}
=======
\documentclass[a4paper]{report}
>>>>>>> 88f87b3... init
=======
\documentclass[a4paper, 12pt]{report}
>>>>>>> f6baf06... add intro

%====================== PACKAGES ======================

\usepackage[french]{babel}
\usepackage[utf8x]{inputenc}
%pour gérer les positionnement d'images
\usepackage{float}
\usepackage{amsmath}
\usepackage{graphicx}
\usepackage[colorinlistoftodos]{todonotes}
\usepackage{url}
%pour les informations sur un document compilé en PDF et les liens externes / internes
\usepackage{hyperref}
%pour la mise en page des tableaux
\usepackage{array}
\usepackage{tabularx}
%pour utiliser \floatbarrier
%\usepackage{placeins}
%\usepackage{floatrow}
%espacement entre les lignes
\usepackage{setspace}
\usepackage{lscape}
%modifier la mise en page de l'abstract
\usepackage{abstract}
%police et mise en page (marges) du document
\usepackage[T1]{fontenc}
\usepackage[top=2cm, bottom=2cm, left=2cm, right=2cm]{geometry}
%Pour les galerie d'images
\usepackage{subfig}

%====================== INFORMATION ET REGLES ======================

\setcounter{secnumdepth}{3}
\setcounter{tocdepth}{3}

\hypersetup{                            % Information sur le document
<<<<<<< HEAD
<<<<<<< HEAD
=======
>>>>>>> f6baf06... add intro
pdfauthor = {Amaury Lavielle},          % Auteurs
pdftitle = {Open Orchestra -
            CMS Symfony2},           % Titre du document
pdfsubject = {Rapport de stage},       % Sujet
pdfkeywords = {Symfony2, CMS},  % Mots-clefs
<<<<<<< HEAD
=======
pdfauthor = {Amaury Lavielle,
            Antoine Lelaisant},          % Auteurs
pdftitle = {Sydonie -
            Création d'un bundle Symfony2 gérant des documents selon la norme FRBR},           % Titre du document
pdfsubject = {Rapport de projet},       % Sujet
pdfkeywords = {FRBR, Sydonie, Document, Gestion de documents},  % Mots-clefs
>>>>>>> 88f87b3... init
=======
>>>>>>> f6baf06... add intro
pdfstartview={FitH}}                    % ajuste la page à la largueur de l'écran

%======================== DEBUT DU DOCUMENT ========================

\begin{document}

%régler l'espacement entre les lignes
\newcommand{\HRule}{\rule{\linewidth}{0.5mm}}

%page de garde
%%%%%%%%%%%%%%%%%%%%%%%%%%%%%%%%%%%%%%%%%
% University Assignment Title Page
% LaTeX Template
% Version 1.0 (27/12/12)
%
% This template has been downloaded from:
% http://www.LaTeXTemplates.com
%
% Original author:
% WikiBooks (http://en.wikibooks.org/wiki/LaTeX/Title_Creation)
%
% License:
% CC BY-NC-SA 3.0 (http://creativecommons.org/licenses/by-nc-sa/3.0/)
%
% Instructions for using this template:
% This title page is capable of being compiled as is. This is not useful for
% including it in another document. To do this, you have two options:
%
% 1) Copy/paste everything between \begin{document} and \end{document}
% starting at \begin{titlepage} and paste this into another LaTeX file where you
% want your title page.
% OR
% 2) Remove everything outside the \begin{titlepage} and \end{titlepage} and
% move this file to the same directory as the LaTeX file you wish to add it to.
% Then add \input{./title_page_1.tex} to your LaTeX file where you want your
% title page.
%
%%%%%%%%%%%%%%%%%%%%%%%%%%%%%%%%%%%%%%%%%

%----------------------------------------------------------------------------------------
%	PACKAGES AND OTHER DOCUMENT CONFIGURATIONS
%----------------------------------------------------------------------------------------
\begin{titlepage}

% \newcommand{\HRule}{\rule{\linewidth}{0.5mm}} % Defines a new command for the horizontal lines, change thickness here

\center % Center everything on the page

%----------------------------------------------------------------------------------------
%	HEADING SECTIONS
%----------------------------------------------------------------------------------------

\textsc{\LARGE Université de Caen Basse-Normandie}\\[1.5cm] % Name of your university/college

\textsc{\Large Rapport de stage}\\[0.5cm] % Major heading such as course name
\textsc{1 avril 2015 - 30 septembre 2015}\\[0.5cm] % Minor heading such as course title


%----------------------------------------------------------------------------------------
%	TITLE SECTION
%----------------------------------------------------------------------------------------

\HRule \\[0.4cm]

{ \huge \bfseries Open Orchestra}\\[0.4cm] % Title of your document

\HRule \\[1.5cm]

%----------------------------------------------------------------------------------------
%	AUTHOR SECTION
%----------------------------------------------------------------------------------------

\begin{minipage}{0.4\textwidth}
\begin{flushleft} \large
\emph{Auteur:}\\

Amaury \textsc{Lavieille}

\end{flushleft}
\end{minipage}
~
\begin{minipage}{0.4\textwidth}
\begin{flushright} \large
\emph{Enseignants:} \\
Marc \textsc{Spaniol}\\
Jean-Marc \textsc{Lecarpentier}
\end{flushright}
\end{minipage}\\[4cm]

%----------------------------------------------------------------------------------------
%	DATE SECTION
%----------------------------------------------------------------------------------------

{\large \today}\\[3cm] % Date, change the \today to a set date if you want to be precise

%----------------------------------------------------------------------------------------
%	LOGO SECTION
%----------------------------------------------------------------------------------------

% Include a department/university logo - this will require the graphicx package

\begin{figure}[H]
    \begin{center}
      \includegraphics[scale=1]{images/LogoUnicaen}
    \end{center}
    \label{frbr-resume}
\end{figure}

%----------------------------------------------------------------------------------------

\vfill % Fill the rest of the page with whitespace

\end{titlepage}

%page blanche
~
%ne pas numéroter cette page
\thispagestyle{empty}

<<<<<<< HEAD
<<<<<<< HEAD
% \renewcommand{\abstractname}{}
\begin{abstract}
Ce projet a été réalisé en binôme par Antoine \textsc{Lelaisant} et Amaury \textsc{Lavieille}. Après de nombreux moments de réflexion et de recherche en commun pour poser les bases du projet, notamment le modèle de document proposé par Sydonie (que nous détaillerons dans ce rapport), nous nous sommes réparti les différentes implémentations à effectuer.

\paragraph{}
Le générateur d'entités et des opérations CRUD\footnote{Acronyme signifiant Create, Read, Update, Delete qui désigne les opérations de base pour la persistance des entités} pour Antoine \textsc{Lelaisant} et l'implémentation du modèle et la négociation de contenu pour Amaury \textsc{Lavieille}.

\paragraph{}
Enfin, nous avons décidé d'expérimenter le travail en \og{}binômage\fg{} pour développer la partie la plus complexe du projet, à savoir, la gestion des documents composites.

\end{abstract}

\chapter*{Remerciements}
\addcontentsline{toc}{chapter}{Remerciements}
Tout d'abord, je tiens à remercier l'ensemble de l'agence Interakting pour m'avoir offert l'opportunité de réaliser mon stage au sein de l'agence.
\paragraph{}
Je tiens aussi à remercier plus particulièrement les différents membres de l'équipe d'Open Orchestra présent durant mon stage, pour leur accueil au sein de l'équipe et leurs conseils

\begin{itemize}
\item[]
\item Nicolas Bouquet, Product Owner
\item[]
\item Nicolas Thal, Scrum master
\item[]
\item Nicolas Anne 
\item[]
\item Noël Gilain
\item[]
\item Romain Boëssel
\item[]
\end{itemize}
=======
 \renewcommand{\abstractname}{}
\begin{abstract}
Ce projet a été réalisé en binôme par Antoine \textsc{Lelaisant} et Amaury \textsc{Lavieille}. Après de nombreux moments de réflexion et de recherche en commun pour poser les bases du projet, notamment le modèle de document proposé par Sydonie (que nous détaillerons dans ce rapport), nous nous sommes réparti les différentes implémentations à effectuer.

\paragraph{}
Le générateur d'entités et des opérations CRUD\footnote{Acronyme signifiant Create, Read, Update, Delete qui désigne les opérations de base pour la persistance des entités} pour Antoine \textsc{Lelaisant} et l'implémentation du modèle et la négociation de contenu pour Amaury \textsc{Lavieille}.

\paragraph{}
Enfin, nous avons décidé d'expérimenter le travail en \og{}binômage\fg{} pour développer la partie la plus complexe du projet, à savoir, la gestion des documents composites.

\end{abstract}


>>>>>>> 88f87b3... init
=======
% \renewcommand{\abstractname}{}
\begin{abstract}
Ce projet a été réalisé en binôme par Antoine \textsc{Lelaisant} et Amaury \textsc{Lavieille}. Après de nombreux moments de réflexion et de recherche en commun pour poser les bases du projet, notamment le modèle de document proposé par Sydonie (que nous détaillerons dans ce rapport), nous nous sommes réparti les différentes implémentations à effectuer.

\paragraph{}
Le générateur d'entités et des opérations CRUD\footnote{Acronyme signifiant Create, Read, Update, Delete qui désigne les opérations de base pour la persistance des entités} pour Antoine \textsc{Lelaisant} et l'implémentation du modèle et la négociation de contenu pour Amaury \textsc{Lavieille}.

\paragraph{}
Enfin, nous avons décidé d'expérimenter le travail en \og{}binômage\fg{} pour développer la partie la plus complexe du projet, à savoir, la gestion des documents composites.

\end{abstract}

\chapter*{Remerciements}
\addcontentsline{toc}{chapter}{Remerciements}
Tout d'abord, je tiens à remercier l'ensemble de l'agence Interakting pour m'avoir offert l'opportunité de réaliser mon stage au sein de l'agence.
\paragraph{}
Je tiens aussi à remercier plus particulièrement les différents membres de l'équipe d'Open Orchestra présent durant mon stage, pour leur accueil au sein de l'équipe et leurs conseils

\begin{itemize}
\item[]
\item Nicolas Bouquet, Product Owner
\item[]
\item Nicolas Thal, Scrum master
\item[]
\item Nicolas Anne 
\item[]
\item Noël Gilain
\item[]
\item Romain Boëssel
\item[]
\end{itemize}
>>>>>>> f6baf06... add intro
\tableofcontents
\thispagestyle{empty}
\setcounter{page}{0}
%ne pas numéroter le sommaire

%espacement entre les lignes d'un tableau
\renewcommand{\arraystretch}{1.5}

%====================== INCLUSION DES PARTIES ======================

~
\thispagestyle{empty}
%recommencer la numérotation des pages à "1"
\setcounter{page}{0}
\newpage
<<<<<<< HEAD
<<<<<<< HEAD
=======
>>>>>>> f6baf06... add intro
\chapter*{Introduction}
\addcontentsline{toc}{chapter}{Introduction}
        \paragraph{}
         Dans le cadre de ma deuxième année du master DNR2I (Master document numérique en réseau, ingénierie de l'Internet), j'ai effectué un stage chez Business \& Decision au sein de l'agence Interakting à Caen.
         \paragraph{}
          Pour une durée de six mois, du 1\up{er} Avril au 30 Septembre 2015, j'ai été intégré au sein de l'équipe de développement du projet Open Orchestra.
         \paragraph{}
         Dans un premier temps, je vais vous présenter succinctement l'entreprise. Ensuite, je vous décrirais le projet Open Orchestra, son origine, ses particularités, etc. Enfin, je vous parlerais de l'organisation et des méthodes mises en place pour la réalisation du projet.

\part{L'entreprise}
\chapter{Présentation de l'entreprise}
\section{Business \& Decision}
        \paragraph{}
        Business \& Decision est un groupe international spécialisé dans trois grand domaines, la Business Intelligence\footnote{La Business Intelligence ou l'informatique décisionnelle représente les différentes solutions informatiques qui permettent l'exploitation des données du entreprise dans le but de faciliter la prise de décision} (BI), la gestion de la relation client et enfin le e-business. 
        \paragraph{}
        Le groupe a était créé en 1992 par Patrick Bensabat, aujourd'hui il est présent dans 15 pays et emplois plus de 2500 personnes.
        \paragraph{}
       	  Événement majeurs et évolution du groupe depuis 1992 : 
       	  \begin{description}
       	  \item[1992] : Création de Business \& Decision par Patrick Bensabat autour de projets de Business Intelligence
       	  \item[1993-1996] : Mise en place des premiers Datawarehouses\footnote{Les Datawarehouses ou entrepôt de données désigne une base de données utilisée pour stocker et structurer des données de production d'une entreprise afin de fournir des informations stratégiques pour la prise de décision} et applications de prévisions financières
       	  \item[1999-2000] : Création de la division CRM (Gestion de la relation client)
       	  \item[2002-2003] : Acquisition en Grande-Bretagne et au Benelux. Création d'agence régionales en France
       	  \item[2004] : Acquisition en Suisse et aux Pays-bas et implantation en Tunisie
       	  \item[2005] : Implantation aux Etats-Unis et acquisition de Metaphora \footnote{Métaphora est une société de service spécialisée dans la conduite du changement. Elle intervient dans toutes les étapes d’un projet pour faciliter l’appropriation du futur système d’informations par les utilisateurs finaux.}
       	  \item[2007-2008] : Création de la marque Interkating
       	  \item[2011] : Déploiement d'offres Cloud et Mobilité. Inauguration du Datacenter éco-responsable d'Eolas
       	  \item[2013] : Lancement de Datalyse, programme de recherche en Big Data. Implantation au Pérou
       	  \item[2014] : Création de Herewecan, agence de communication digitale, Lancement des pôles d'expertise Big Data, Transformation et Hub Mobile )
       	  \end{description}
       	  
\section{Interakting}

        \paragraph{}
        Interakting est l'agence web du groupe B\&D  composé de 340 employés répartis principalement entre Caen et Paris, les équipes de l'agence interviennent sur des projets web  de grande envergure (Canal +, Bnp Paribas, Inra, Moët Hennessy, PSA Peugeot Citroën, ...).
        
         \paragraph{}
        Les différents projets de l'agence sont développés avec différent outils : eZ publish (système de gestion de contenu créé par l'entreprise eZ Systems As), PHP Factory (plateforme réalisé par conjointement par Interakting et Zend Technologies) et de plus en plus avec Symfony2 (framework php écrit par SensionLabs) avec notamment Open Orchestra.

        
        
        
\part{Open Orchestra}
\part{Organisation du développement}
<<<<<<< HEAD
=======

\part{Présentation}
    \chapter*{Introduction}
\addcontentsline{toc}{chapter}{Introduction}
        \paragraph{}
         Dans le cadre de ma deuxième année du master DNR2I (Master document numérique en réseau, ingénierie de l'Internet), j'ai effectué un stage chez Business \& Decision au sein de l'agence Interakting à Caen.
         \paragraph{}
          Pour une durée de six mois, du 1\up{er} Avril au 30 Septembre 2015, j'ai été intégré au sein de l'équipe de développement du projet Open Orchestra.
         \paragraph{}
         Dans un premier temps, je vais vous présenter succinctement l'entreprise. Ensuite, je vous décrirais le projet Open Orchestra, son origine, ses particularités, etc. Enfin, je vous parlerais de l'organisation et des méthodes mises en place pour la réalisation du projet.


    \input{./sydonie.tex}

    \input{./symfony2.tex}

\part{Implémentation}
    \input{./model-implementation.tex}

    \input{./generator-entity.tex}

    \input{./generator-crud.tex}

    \input{./content-negociation.tex}

    \input{./document-composite.tex}
\part{Conclusions et perspectives}
    \chapter*{Bilan}
\addcontentsline{toc}{chapter}{Bilan}
Durant le stage, ce fut un réel plaisir de travailler avec l'équipe du projet Open Orchestra et plus en général avec les différentes personnes d'Interakting que j'ai côtoyé à l'agence de Caen ou de Paris lors de mes déplacements.
\paragraph{}
J'ai été directement intégré au projet comme aurait pu l'être un autre développeur de l'agence ce qui a été très constructif pour moi et important de voir que l'on me faisait confiance.
\paragraph{}
Travailler sur un projet Open Source d'un point de vue plus orienté éditeur que intégrateur (réalisation de site pour des clients) a été très intéressants et motivant.  Cela ma permit de voir et mettre en pratique la méthode agile Scrum sur un projet concret.
\paragraph{}
D'un point vue de technique, cela m'a permis de comprendre et approfondir des concepts avancés de  Symfony2  mais pas uniquement. Plus généralement, j'ai pu apprendre énormément sur les bonnes pratiques (test, psr2, principe solid @TODO solid
) mais aussi de découvrir et d'utiliser de nouveau outil/langages (Travis, Ansible, Vagrant, CoffeeScript, Jmeter). 
De plus, durant mon stage j'ai assisté à des formations sur différent sujet (Git, Travis, Scrutinizer) proposé tous les jeudis de 13h à 15h par des développeurs de l'entreprise.
\paragraph{}
Pour conclure, ce stage à largement répondu à mes attentes que ce soit sur le projet en lui-même ou sur le point technique.

>>>>>>> 88f87b3... init
=======
>>>>>>> f6baf06... add intro
\input{./appendix.tex}

\newpage

%récupérer les citation avec "/footnotemark"
\nocite{*}

%choix du style de la biblio
\bibliographystyle{plain}
%inclusion de la biblio
\bibliography{bibliographie.bib}
%voir wiki pour plus d'information sur la syntaxe des entrées d'une bibliographie

\end{document}